\documentclass{article} 

\usepackage{mdframed}
\usepackage{listings}
\usepackage{hyperref}
\usepackage{amsfonts}
\usepackage{amsmath}

\title{CSE160 Homework 1}
\author{Elijah Hantman}

\renewcommand{\familydefault}{\sfdefault}
%\renewcommand{\spn}[1]{\langle #1 \rangle}


\begin{document}
    \maketitle

    {\large Part 1}
    \begin{enumerate}
        \item Normalize vector $\vec{v} = \begin{bmatrix}
                1\\
                2\\
                3
            \end{bmatrix} $. Round to 2 decimal places.

            \begin{mdframed}
                \begin{gather}
                    ||\vec{v}|| = \sqrt{1 + 2^2 + 3^2}\\
                    ||\vec{v}|| = \sqrt{1 + 4 + 9} = \sqrt{14}\\
                    \frac{\vec{v}}{||\vec{v}||} = \begin{bmatrix}
                        0.27\\
                        0.53\\
                        0.80
                    \end{bmatrix}
                \end{gather}
            \end{mdframed}
        \item Simplify and Find the Length of the Vector
            \begin{displaymath}
                3 * [1,1] + 2 * [-1, 1]
            \end{displaymath}
            \begin{mdframed}
                \begin{gather}
                    3 * [1,1] + 2 * [-1, 1]\\
                    = [3, 3] + [-2, 2]\\
                    = [1, 5]
                \end{gather}

                And the length:
                \begin{gather}
                    |[1, 5]| = \sqrt{1 + 5^2}\\
                    = \sqrt{26}\\
                    \approx 5.099
                \end{gather}
            \end{mdframed}
            \break
        \item Calculate the Cross Product
            \begin{enumerate}
                \item $[0,1,1] \times [1,1,0]$
                    \begin{mdframed}
                        \begin{gather}
                            \det{\begin{bmatrix}
                                i & j & k\\ 
                                0 & 1 & 1\\
                                1 & 1 & 0
                        \end{bmatrix}} = \\
                        i(1)(0) + j(1)(1) + k(0)(1) - i(1)(1)
                        - j(0)(0) - k(1)(1) = \\
                        j - i - k = j - k = [-1, 1, -1]
                        \end{gather}       
                    \end{mdframed}
                \item $[2, 3, 4] \times [1, 0, 0]$
                    \begin{mdframed}
                        \begin{gather}
                            \det{
                                \begin{bmatrix}
                                    i & j & k\\ 
                                    2 & 3 & 4\\
                                    1 & 0 & 0
                                \end{bmatrix}
                            } \\
                            = i(3)(0) + j(4)(1) + k(2)(0)
                            - i(4)(0) - j(2)(0) - k(3)(1)\\
                            = 4j - 3k = [0, 4, -3]
                        \end{gather}
                    \end{mdframed}
                \item $[0, 3, 4] \times [2, 2, 2]$
                    \begin{mdframed}
                        \begin{gather}
                            \det{
                                \begin{bmatrix}
                                    i & j & k\\ 
                                    0 & 3 & 4\\
                                    2 & 2 & 2
                                \end{bmatrix}
                            } = \\
                            i(3)(2) + j(4)(2) + k(0)(2)
                            - i(4)(2) - j(0)(2) - k(3)(2)\\
                            = 6i + 8j - 8i - 6k\\
                            = -2i + 8j -6k = [-2, 8, -6]
                        \end{gather}
                    \end{mdframed}
            \end{enumerate}
            \break
        \item Calculate the Dot Product
            \begin{enumerate}
                \item $[1, 0, 1] \cdot [0, 1, 1]$
                    \begin{mdframed}
                        \begin{gather}
                            1 \cdot 0 + 0 \cdot 1 + 1 \cdot 1\\
                            = 1
                        \end{gather}
                    \end{mdframed}
                \item $[0, 3, 4] \cdot [1, 0, 0]$
                    \begin{mdframed}
                        \begin{gather}
                            0 \cdot 1 + 3 \cdot 0 + 4 \cdot 0\\
                            = 0
                        \end{gather}
                    \end{mdframed}
                \item $[2, 3, 4] \cdot [6, 4, 3]$
                    \begin{mdframed}
                        \begin{gather}
                            2 \cdot 6 +
                            3 \cdot 4 +
                            4 \cdot 3\\
                            = 12 + 12 + 12 = 36
                        \end{gather}
                    \end{mdframed}
            \end{enumerate}
            \break
        \item Consider a triangle formed by connecting the three points $p_1 = (0, 0, 0)$, $p_2 = (1, 0, 0)$
            and $p_3 = (1, 1, 1)$.
            \begin{mdframed}
                First we find $p_2 - p_1$ and $p_3 - p_1$

                \begin{gather}
                    p_2 - p_1 = (1 - 0, 0 - 0, 0 - 0) = (1, 0, 0)\\
                    p_3 - p_1 = (1 - 0, 1 - 0, 1 - 0) = (1, 1, 1)
                \end{gather}

                Here, since $p_1 = \vec{0}$ we don't technically have to
                subtract, but a general method would subtract

                Then we use the cross product.

                \begin{gather}
                    p_2 \times p_3 = (0 * 1 - 0 * 1, 0 * 1 - 1 * 1, 1 * 1 - 0 * 1)\\
                    = (0, -1, 1)
                \end{gather}

                Now, the cross product is also known to be perpendicular
                to its factors, meaning that $(0, -1, 1)$ has a positive
                z direction and is perpendicular to the triangle.

                To find the area we simply divide the length of the cross product
                by 2 since the cross product is the area of a parallelogram.

                \begin{gather}
                    |(0, -1, 1)| = \sqrt{1 + 1} = \sqrt{2}\\
                    A = \frac{\sqrt{2}}{2}
                \end{gather}
            \end{mdframed}

        \item Calculate the Matrix 
            \begin{enumerate}
                \item 
                    \begin{gather}
                        \begin{bmatrix}
                            1 & 2 & 5\\ 
                            -1 & -1 & 1\\
                            4 & 4 & -2
                        \end{bmatrix} \cdot
                        \begin{bmatrix}
                            1\\ 
                            2\\
                            3
                        \end{bmatrix}\\
                        = \begin{bmatrix}
                            1\\ 
                            -1\\
                            4
                        \end{bmatrix} +
                        \begin{bmatrix}
                            4\\ 
                            -2\\
                            8
                        \end{bmatrix} +
                        \begin{bmatrix}
                            15\\ 
                            3\\
                            -6
                        \end{bmatrix}\\
                        = \begin{bmatrix}
                            20\\ 
                            0\\
                            6
                        \end{bmatrix}
                    \end{gather}
                \item 
                    \begin{gather}
                        \begin{bmatrix}
                            2 & 0 & 3\\ 
                            0 & -1 & 2\\
                            3 & 2 & -2
                        \end{bmatrix} \cdot
                        \begin{bmatrix}
                            -1 & -4 & 1\\ 
                            1 & -1 & 4\\
                            0 & 0 & 5
                        \end{bmatrix}\\
                        =
                        \begin{bmatrix}
                            -2 & -8 & 17\\ 
                            -1 & 1 & 6\\
                            -1 & -14 & 1
                        \end{bmatrix}
                    \end{gather}
            \end{enumerate}
        \item Consider Two lines $y = 4/3 x - 1$ and
            $y = 0$
            \begin{enumerate}
                \item Find the intersection of the two lines

                    \begin{mdframed}
                        \begin{gather}
                            0 = 4/3 x - 1\\
                            1 = 4/3 x\\
                            3/4 = x\\
                            y = 4/3(3/4) - 1\\
                            y = 1 - 1\\
                            y = 0
                        \end{gather}

                        The intersection is $(0.75, 0)$.
                    \end{mdframed}
                \item Are these lines Perpendicular
                    \begin{mdframed}
                        No.

                        We can directly calculate this by taking
                        the dot product of the direction of each line.
                        This works since $cos(90^{\circ}) = 0$. If
                        the result is nonzero then the lines must not
                        be perpendicular.

                        To get the directions we take rise over run.

                        \begin{gather}
                            d_0 = (3, 4)\\
                            d_1 = (1, 0)\\
                            d_0 \cdot d_1 = (3 * 1 + 4 * 0) = 3
                        \end{gather}
                    \end{mdframed}
            \end{enumerate}
    \end{enumerate}
    \break

    {\large Part 2}

    \begin{enumerate}
        \item Define OpenGL, OpenGL ES and WebGL. Describe their relationship.
            \begin{mdframed}
                OpenGL is a standard API for communicating to graphics
                drivers and their associated graphics processing units.
                The main characteristics of OpenGL is being fairly high level
                in terms of API and being an open standard meaning it is not
                propriatary to any specific platform or device.

                OpenGL ES is a version of the OpenGL standard geared towards
                mobile devices and embedded systems, hence the name
                OpenGL ES (Embedded Systems).

                WebGL is an API provided by browsers to javascript websites
                in order for web applications to utilize GPU acceleration.
                It is based on OpenGL ES.

                Each of these are related in the following ways. OpenGL is
                the oldest. OpenGL ES was based on an OpenGL version,
                meaning OpenGL ES 1.1 was based on OpenGL 1.5, and OpenGL ES
                2.0 was based on OpenGL 2.0. WebGL is based on OpenGL ES,
                and tracks along with the ES standard.
            \end{mdframed} 
        \item What is the difference in the software architecture of a webpage and a webpage using WebGL.
            Describe all components involved in the two cases.
            \begin{mdframed}
                A standard webpage uses 3 languages. First html
                defines the main structure of the page, defining
                each element.

                Then CSS is used to layout the page, setting colors,
                styling and how html elements should be organized.

                Javascript is then run in an event loop which provides
                interactivity in the browser without needing a full
                roundtrip request.

                A WebGL webpage simply adds a connection to OpenGL ES
                which provides an api in javascript.

                These translate to the following software browser components,
                an html rendering engine that handles CSS,
                a javascript runtime,
                and an OpenGL ES context, which is provided by the
                driver.
            \end{mdframed}
        \item In HTML, what is the <canvas> tag?
            \begin{mdframed}
                The Canvas tag provides a buffer and html element
                which javascript can render to. It can be written
                to via cpu side instructions, and in the context
                of WebGL it is a buffer which can be rendered
                to and displayed on the webpage.
            \end{mdframed}
        \item In Javascript, what line of code one writes to retrieve a 2D rendering context?
            \begin{mdframed}
                Assuming you already have the canvas, meaning you did something like

                \begin{lstlisting}[language=C]
var canvas = document.getElementbyID("name");
                \end{lstlisting}

                According to the Mozilla developer documentation
                to get a webgl context you do:

                \begin{lstlisting}[language=C]
var gl = canvas.getContext("webgl"); 
                \end{lstlisting}

                According to the textbook the specific string
                used to get the context varies between browsers
                so you would normally replace it with a function.
            \end{mdframed}
    \end{enumerate}


\end{document}
